% Generated draft: requires mathematical and source review.
\documentclass[11pt]{article}
\usepackage[margin=25mm]{geometry}
\usepackage{fontspec}
\setmainfont{Latin Modern Roman}
\usepackage{amsmath,amssymb,mathtools,unicode-math}
\setmathfont{Latin Modern Math}
\usepackage{enumitem,needspace,xurl,hyperref}
\setlength{\parindent}{0pt}
\setlength{\parskip}{5pt}
\setlength{\emergencystretch}{3em}
\newcommand{\N}{\mathbb{N}}
\newcommand{\Z}{\mathbb{Z}}
\newcommand{\Q}{\mathbb{Q}}
\newcommand{\R}{\mathbb{R}}
\newcommand{\C}{\mathbb{C}}
\newcommand{\F}{\mathbb{F}}
\newcommand{\A}{\mathbb{A}}
\newcommand{\PP}{\mathbb{P}}
\newcommand{\E}{\mathbb{E}}
\newcommand{\eps}{\varepsilon}
\newcommand{\dd}{\,\mathrm{d}}
\begin{document}
\section*{UnsolvedMath ID: 801017}
\textbf{Source problem code:} catalogue entry
\subsection{Definitions and mathematical statement}
The supplied record contains two questions referring to a preceding Velickovic--Woodin theorem whose statement is absent. Consequently, the target theorem cannot be reconstructed from the supplied evidence.

Under the customary interpretation of the notation, \(\Sigma^1_1(y)\) denotes the class of sets of reals that are effectively analytic relative to the real oracle \(y\), and \(\omega_1^{CK(z)}\) denotes the least ordinal having no \(z\)-computable well-order presentation. The source does not explicitly specify the ambient space of reals or its coding.

\begin{enumerate}
\item Is there a recursion-theoretic proof of the Velickovic--Woodin theorem referred to in the source as the ``above theorem''? The theorem and the intended requirements on such a proof are not supplied.
\item For a parameter \(y\) and every \(\Sigma^1_1(y)\) set \(A\), does
\[
\exists x\in A\;\bigl(\omega_1^{CK(x)}>\omega_1^{CK(y)}\bigr)
\]
imply that the Velickovic--Woodin theorem holds for \(A\)? Here ``the theorem holds for \(A\)'' retains the source's unresolved reference; its conclusion and any additional hypotheses are unavailable. The scope and domain of the parameter \(y\) are implicit in the source.
\end{enumerate}
These are incomplete formulation targets pending recovery of the referenced theorem, not fully specified mathematical assertions.
\subsection{Short English statement}
The source asks for a recursion-theoretic proof of an unspecified Velickovic--Woodin theorem and asks whether that theorem extends to every effectively analytic set relative to \(y\) containing an \(x\) with \(\omega_1^{CK(x)}>\omega_1^{CK(y)}\). The missing theorem prevents a complete precise statement.
\subsection{Sources}
\begin{itemize}
\item Local UnsolvedMath registry: permanent ID \(801017\), source code \texttt{catalogue entry}, title ``Open Problems in Mathematical Logic --- Problem 29.'' The registry attributes the record to Barmpalias, Gao et al., with proposed year 2026.
\item Supplied source reference: Barmpalias, Gao et al., \emph{Open Problems in Mathematical Logic} (2026), Problem 29, \url{https://arxiv.org/abs/2608.26628}. The record reports access on 2026-10-16 and extraction by \texttt{arxiv\_latex}; these are supplied provenance claims, not independently verified access or source verification.
\item The source names Velickovic--Woodin in connection with the missing preceding theorem but supplies neither that theorem nor a separate bibliographic citation for it. Upstream revision: \texttt{d371ee917ac9647fbb06cd6fc529deec6f3dda82}. The registry labels the problem ``open''; this is supplied status metadata, not an independently verified mathematical assessment.
\end{itemize}
\end{document}
